%============================================================ BLOQUE 7 · Conclusiones (10 min)
\ugrsection{07}{Conclusiones y preguntas}{Inicio: 1:50 · 10 min\\Lo que conviene recordar}{Tres mensajes clave}

\begin{frame}{Para llevarse a casa}{Tres mensajes}
  \note{Infraestructura y pasarela; memoria y validación; agente y permisos. NEXT SLIDE: Límites problemas}
  \begin{enumerate}
    \item Trata el LLM como \textbf{infraestructura}: controla el camino de los datos, estandariza en la API compatible con OpenAI y pon una \textbf{pasarela} delante.
    \item La \textbf{memoria y la concurrencia} deciden qué puedes servir. \textbf{Restringe las salidas, valida y monitoriza}; empieza en pequeño y escala en horizontal.
    \item Un agente es un \textbf{bucle con herramientas}: MCP y las skills lo amplían; los permisos y el sandbox lo contienen.
  \end{enumerate}
  \keymsg{\textbf{Mensaje final:} el valor no está solo en el modelo, sino en la infraestructura y las prácticas que lo rodean.}
\end{frame}

\begin{frame}{Límites y problemas abiertos}{Dónde estamos y qué queda por hacer}
  \note{Alucinaciones, modelos pequeños, coste GPU, mantenimiento y evaluación. NEXT SLIDE: Prueba kit}
  \begin{columns}[T]
    \begin{column}{0.5\textwidth}
      \themeblock{nBrown}{Modelos}{%
        Alucinaciones y conocimiento desactualizado.\\
        Los modelos abiertos van por detrás de la frontera en razonamiento difícil.\\
        Los modelos pequeños fallan con herramientas en tareas largas.}
    \end{column}
    \begin{column}{0.5\textwidth}
      \themeblock{nD}{Operación}{%
        Coste de GPU, energía y adquisición.\\
        Mantenimiento constante: modelos, motores, parches de seguridad.\\
        La evaluación y el cumplimiento del AI Act exigen un esfuerzo real.}
    \end{column}
  \end{columns}
\end{frame}

\begin{frame}{Pruébalo tú mismo}{Ejercicios con el kit de demo y dónde seguir leyendo}
  \note{Invitar a probar herramientas, modelos, esquemas y skills. NEXT SLIDE: Pregunta sobre MiIA}
  \begin{columns}[T]
    \begin{column}{0.56\textwidth}
      \small
      \begin{enumerate}
        \item Añade una herramienta a \texttt{mcp/server.py}: el agente la descubre solo.
        \item Registra un segundo modelo en LiteLLM.
        \item Cambia el esquema de \texttt{03} para otro documento.
        \item Escribe tu propia skill en \texttt{workspace/skills/}.
        \item Prueba un modelo más grande con las herramientas.
      \end{enumerate}
    \end{column}
    \begin{column}{0.42\textwidth}
      \themeblock{rojoUGR}{Documentación}{%
        \texttt{docs.vllm.ai}\\
        \texttt{docs.litellm.ai}\\
        \texttt{modelcontextprotocol.io}\\
        \texttt{docs.openwebui.com}}
    \end{column}
  \end{columns}
\end{frame}
